\section{结论}

本文构建了一套面向综合 AI 应用机会的多源智能挖掘与价值评估系统，并在 B01 公开数据批次上完成了从采集、清洗、机会判别、代理复核、主题聚合到价值排序的可复现流程。系统保留了来源、正文、质量标记、模型证据、代理裁定和主题成员之间的处理关系，能够区分原模型标签、盲复核标签和 AI 代理标签。

B01 最终清洗得到 2,561 条记录。原模型覆盖全量记录并标记 269 条候选机会；509 条分层复核队列经过两轮辅助复核和冲突裁定后，形成 100 条 AI 代理机会候选。基于代理候选的语义聚合得到 14 个探索性主题，但其中 12 个为单记录主题，轮廓系数为 0.204。第五项任务已经实现复核工作台和决策者主题优先级视图，C2 策划案仍属于待确认框架。结果说明系统已经能够把复杂公开文本组织为可审计的候选链路，也同时暴露出来源偏差、正文缺失、定向抽样、模型相关性和主题稀疏等问题。

本文的主要结论是：在缺少真人真值时，完整记录处理链路和显式不确定性比输出一个看似精确的热点排名更重要。AI 代理结果可以用于系统调试、候选筛选和论文中的探索性案例，但不能替代人工真值，也不能支持准确率、召回率和外部市场占比等强结论。后续完成真人复核、分层随机评价和主题边界确认后，系统才能形成正式比赛结果和可执行的策划案。

因此，B01 的价值不在于证明某个 AI 方向已具备投资或上线条件，而在于展示从公开记录到决策候选的可核验处理过程。后续验证应先对全部原模型机会记录开展真人确认，再在非机会记录中按来源和正文完整度分层随机抽样；完成标签一致性检查后，重新计算主题、来源覆盖、价值分和排名稳定性。策划案的用户范围、功能优先级与风险控制也需由成员级原文和后续访谈确认。完成这些工作之前，C1 和 C2 仅能用于提出问题与安排验证，不能代表实际市场规模。

后续结果应继续保留批次、来源和标签版本，使人工确认后的正式结果能够与当前探索性结果逐项区分。
